\documentclass[11pt]{article}
\usepackage[a4paper, margin=2.5cm]{geometry}

\usepackage{libertine}             % nice font
\usepackage{libertinust1math}      % nice math font
\usepackage[utf8]{inputenc}        % allow utf-8 input
\usepackage[T1]{fontenc}           % use 8-bit T1 fonts
\usepackage{microtype}             % microtypography

% https://tex.stackexchange.com/questions/50747/options-for-appearance-of-links-in-hyperref
% \usepackage[colorlinks,citecolor=magenta]{hyperref}  % hyperlinks
\usepackage{hyperref}
\usepackage{natbib}                % citations
\usepackage{booktabs}              % professional-quality tables
\usepackage{enumerate}             % better lists
\usepackage{lineno}                % line numbers


\usepackage{titling}
\setlength{\droptitle}{-2cm} % move up the title to have more space

\title{\textbf{Project Title}\\
\Large{732A76 Research Project}}

\author{%
    Name
}

\date{\today}

% comment out the following line to hide line numbers
\linenumbers

\begin{document}

\maketitle

\begin{abstract}
An abstract to summarize the problem and your findings.
Obey the page limit of 3-5 pages, excluding appendix and references. Your report must be fully understandable within those 3-5 pages. This page limit includes the title and abstract. Do not include a table of contents. A violation of the page limit will not be accepted and the report has to be resubmitted.
\end{abstract}

\section{Introduction}

A brief introduction to the problem.

\section{Your Contribution}

What did you do? Adjust the section title.

\section{Experiments}

Experimental details.

\textbf{Data.} Asd

\textbf{Baselines.} Sdf

\textbf{Etc.} qwerty

\section{Discussion}

How to interpret the results you obtained?

\section{Conclusion}

Very brief summary and conclusion of your work.

\section*{References}

You can use (and reference) online sources, but use your own words.
\textbf{Note:} Copy and paste from online sources or past reports will be identified in our plagiarism check and will lead to automatic failure and (in serious cases) reporting to the University's disciplinary board.

\appendix

\section{Code}

\paragraph{Code template:}
Your report should include the main code in the appendix section. To make your code readable in the PDF file, you can use a code-style template by adding the lines 
\begin{verbatim}
    \usepackage{xcolor}
    \definecolor{codegreen}{rgb}{0,0.6,0}
    \definecolor{codegray}{rgb}{0.5,0.5,0.5}
    \definecolor{codepurple}{rgb}{0.58,0,0.82}
    \definecolor{backcolour}{rgb}{0.95,0.95,0.92} 
    \usepackage{listings}
    \lstdefinestyle{mystyle}{
        backgroundcolor=\color{backcolour},   
        commentstyle=\color{codegreen},
        keywordstyle=\color{magenta},
        numberstyle=\tiny\color{codegray},
        stringstyle=\color{codepurple},
        basicstyle=\footnotesize\ttfamily,
        breakatwhitespace=false,    
        breaklines=true, captionpos=b,    
        keepspaces=true, numbers=left,   
        numbersep=5pt, showspaces=false, 
        showstringspaces=false,
        showtabs=false, tabsize=2}
    \lstset{style=mystyle}
\end{verbatim}
before \verb|begin{document}| and use \verb|\lstinputlisting[language=python]{code.py}| to insert code to the appendix section.

\textit{Another choice}: If you write your code using jupyter notebook, you can also directly export the notebook to a PDF and concatenate it to the report as an additional appendix. %Here we recommend a useful website for jointing PDF: \url{https://pdfjoiner.com}. 


\section{Style}
This section was copied from advice for writing a masters thesis but is considered also helpful for this report.

\subsection*{Writing}
\begin{enumerate}
    \item \textbf{Write in English.} Either UK or US English is fine, but try to be consistent.
    
    \item \textbf{Write correct English.} Your editor most likely supports spell correction. Use it!
    
    \item \textbf{Use active voice.} ``We trained the neural network'' is preferable over ``The neural network was trained''.  
   
    \item \textbf{Use consistent terminology.} Many concepts are known under different names (e.g. artificial neural networks, neural networks, deep neural networks, deep learning), but to avoid confusion you should pick one and consistently use only that. 
    
    \item \textbf{Avoid abbreviations}. As a rule of thumb, only ever consider abbreviations that are either (i) generally used and understood by everyone in the field, e.g., PDE for Partial Differential Equation, or (ii) used everywhere in your report, e.g., at least three times in the same section. Consider adding a list of abbreviations/acronyms in the beginning.

    \item \textbf{Use correct references.} When referring to a specific algorithm, equation, figure, section, or table, the (usually lowercase) object gets capitalized, e.g., ``as we can see in Figure~5'' vs ``as depicted in the figure''. Use \texttt{\textbackslash eqref} when referencing equations. 

    \item \textbf{Only cite academic work.} If you need a statement from a company website or blog, set a footnote. Only cite proper academic output. Some software packages can be seen as academic work. If you are unsure what to cite, cite the paper that introduced the idea/concept and an article or book describing it best. Do not cite random papers for well-known concepts!
    
    \item \textbf{Use a consistent citation style.} The Harvard style is preferred. Note that there is a difference between an active citation (you would read it out loud) and a passive one (you would not read it out loud). Two examples: Author et al. (2023) show a result. A new result was shown (Author et al., 2023).

\end{enumerate}

\subsection*{Mathematics}
Mathematical writing is fundamentally different from creative and expository writing. In addition to the advice below, please refer to the ``Ten Simple Rules for Mathematical Writing'' by Dimitri P. Bertsekas\footnote{\url{https://www.mit.edu/~dimitrib/Ten_Rules.html}}.
\begin{enumerate}
    \item \textbf{Use \LaTeX.} Try to avoid writing in MS Word whenever possible.
    
    \item \textbf{Never start a sentence with a mathematical symbol.} Never.
    
    \item \textbf{Respect conventions.} Although you are technically free to call your inputs $y$ and your outputs $x$, the reader will dislike you for it.
    
    \item \textbf{Math is always a part of a sentence.} There is no equation between two sentences. An equation is a part of a sentence and thus needs proper punctuation marks!
    
    \item \textbf{Use single-letter variables.} If necessary, use subscripts for additional clarity. It is also advised to specify the space of newly introduced variables, e.g., $\lambda\geq0$ (or $\lambda\in\mathbb{R}_{\geq0}$, or $\lambda\in\mathbb{R}_0^+$). In addition, be consistent. The meaning of a variable should not change during sections.    
\end{enumerate}

\subsection*{Presentation}
\begin{enumerate}
     \item \textbf{Start with the big picture.} Start most sections and subsections by reminding the reader about the big picture, where you are, and what you'll describe next.
    
    \item \textbf{Present things in logical, not chronological, order.} The point of the report is to communicate your findings in the way that makes the most sense for the \emph{reader}. This is often not the same way that you arrived at them. See the note on ``How to write a scientific report''\footnote{\url{https://www.it.uu.se/edu/course/homepage/projektTDB/vt05/examination/rapport.pdf}} by Michael Thuné for a longer explanation of this point.
    
     \item \textbf{Structure your text.} Use sections, subsections, subsubsections, and paragraphs wisely. Not every subsection needs a number, but headings help to structure the text. Thus, if you have two pages of full text, it might be wise to break the text and introduce a heading somewhere.
     
    \item \textbf{Table of contents.} If your table of contents is longer than 1.5-2 pages, reduce. The reduction should not be done by removing structure but rather by not enlisting every subsubsection in the table of contents. Not every subsection needs a number!

    \item \textbf{State your contributions.} It should be clearly stated what your contributions are. If two students jointly write the thesis, it should be clear which student contributed what. Ideally, this is done on a chapter basis or within an \textit{outline} paragraph at the end of the introduction. 

    \item \textbf{Use tables sparingly.} Use \texttt{booktabs} and read \url{https://nhigham.com/2019/11/19/better-latex-tables-with-booktabs/}. Table captions go above tables while figure captions are placed below figures.
    
    \item \textbf{Use images.} Images often help to illustrate concepts and ease understanding. State the image sources if necessary and use appropriate image qualities. 
    
    \item \textbf{Do proper figures.} Figures should have clearly defined axis labels and include a legend. Any text in the figure should approximately match the font size of the thesis main text, thus, it should be readable without zooming in. A transparent grid in the background usually improves the appearance and readability of the figure. Try to use colors that also work for colorblind readers and look good when the thesis is printed in black and white. Images, e.g., created with \texttt{matplotlib}, are best saved in pdf format. 
\end{enumerate}



\end{document}